%!TEX program = xelatex
% ==============================================================================
% FRONTESPIZIO A COLORI: RACCOLTA ESERCITAZIONI MATLAB - INDUSTRIAL ROBOTICS
% ==============================================================================
\documentclass[11pt, a4paper]{article}

% BB Localizzazione Linguistica e Font XeLaTeX
\usepackage{fontspec}
\usepackage[italian]{babel}
\setmainfont{Latin Modern Roman}
\setsansfont{Latin Modern Sans}
\setmonofont{Consolas}[Scale=0.92]

% BB Geometria e Margini della Copertina
\usepackage{geometry}
\geometry{
    a4paper,
    left=2.3cm,
    right=2.3cm,
    top=2.5cm,
    bottom=2.5cm
}

% BB Palette Colori Istituzionale (Coerente con TesinaNEW.tex)
\usepackage[table]{xcolor}
\definecolor{primaryNavy}{RGB}{20, 52, 85}     % Blu Navy istituzionale
\definecolor{accentBlue}{RGB}{41, 128, 185}    % Blu tecnico per dettagli
\definecolor{slateGray}{RGB}{108, 117, 125}    % Grigio sobrio per testi secondari
\definecolor{softBg}{RGB}{248, 249, 250}       % Grigio perlaceo di sfondo
\definecolor{borderGray}{RGB}{222, 226, 230}   % Linee di separazione

% BB Grafica e Inclusione Loghi
\usepackage{graphicx}
\graphicspath{{./Images/}}

% BB Tabelle e Tipografia
\usepackage{booktabs}
\usepackage{array}

% BB Metadati PDF
\usepackage{hyperref}
\hypersetup{
    hidelinks,
    pdfauthor={Francesco Lo Verde},
    pdftitle={Esercitazioni di Laboratorio MATLAB - Industrial Robotics},
    pdfsubject={Raccolta delle esercitazioni di cinematica dei manipolatori robotici},
    pdfkeywords={Industrial Robotics, MATLAB, Cinematica Diretta, Cinematica Inversa, Jacobiano Geometrico}
}

\begin{document}
\pagestyle{empty}

\begin{titlepage}
    \centering
    \sffamily

    % --- SEZIONE SUPERIORE: LOGO E UNIVERSITÀ ---
    \vspace*{0.2cm}
    \includegraphics[width=0.34\textwidth]{LogoUnipg.png}\par
    \vspace{0.8cm}

    {\scshape\Large\bfseries\color{primaryNavy} Università degli Studi di Perugia\par}
    \vspace{0.25cm}
    {\large\bfseries\color{slateGray} Dipartimento di Ingegneria\par}

    \vfill

    % --- SEZIONE CENTRALE: CORSO E TITOLO PRINCIPALE ---
    {\Large\bfseries\color{accentBlue} CORSO DI INDUSTRIAL ROBOTICS\par}
    \vspace{0.6cm}

    \begin{minipage}{0.95\linewidth}
        \centering
        {\fontsize{24}{30}\selectfont \bfseries \color{primaryNavy} Esercitazioni MATLAB\par}
    \end{minipage}

    \vspace{0.7cm}
    {\color{accentBlue}\rule{0.65\linewidth}{1.4pt}}\par
    \vspace{0.7cm}

    {\large\color{slateGray} Rappresentazioni Orientamento, Cinematica Diretta Manipolatori,\\Disaccoppiamento Polso - Struttura Portante, Cinematica Inversa Manipolatore Stanford, Jacobiano Geometrico E Metodi Per Cinematica Inversa\par}

    \vfill

    % --- SEZIONE INFERIORE: DATI DELLO STUDENTE ---
    \begin{minipage}{0.85\linewidth}
        \large
        \begin{tabular}{@{} l l}
            \multicolumn{2}{l}{\textbf{\color{primaryNavy}Studente:}} \\
            \addlinespace[0.4em]
            Nome:  & \textbf{Francesco Lo Verde} \\
            Email: & \texttt{francesco.loverde@studenti.unipg.it} \\
        \end{tabular}
    \end{minipage}

    \vspace{1.5cm}

    % --- PIÈ DI PAGINA: ANNO ACCADEMICO ---
    {\large\bfseries\color{slateGray} Anno Accademico 2025--2026\par}
    \vspace*{0.2cm}

\end{titlepage}

\end{document}

